\begin{solapartat}
\begin{minipage}[c]{0.68\linewidth}
El camp $\vec{B}_A$ (camp magnètic creat pel fil A) és tangent a la circumferència centrada en A i el camp $\vec{B}_B$ (camp magnètic creat pel fil B) és tangent a la circumferència centrada en B.
\end{minipage}\hfill
\begin{minipage}[c]{0.28\linewidth}
\figura[0.85]{sol_fig1.png}
\end{minipage}

\punts{0,4} Esquema
\figura[0.25]{sol_fig2.png}

\punts{0,2}
\[ B_A = B_B = \frac{\mu_0 I}{2\pi d} = \frac{4\pi\times 10^{-7}\times 0,30}{2\pi\times 0,20} = 3,0\times 10^{-7}\ \mathrm{T} \]

\punts{0,2} $B_{total} = \sqrt{B_A^2 + B_B^2}$

\punts{0,2} $B_{total} = \sqrt{2\times\left(3,0\times 10^{-7}\right)^2} = 3,0\times 10^{-7}\sqrt{2} = 4,2\times 10^{-7}\ \mathrm{T}$
\end{solapartat}

\begin{solapartat}
La força $\vec{F}_\mathrm{A}$ és perpendicular a $\vec{B}_\mathrm{A}$ i al corrent $I_\mathrm{c}$.\\
La força $\vec{F}_\mathrm{B}$ és perpendicular a $\vec{B}_\mathrm{B}$ i al corrent $I_\mathrm{c}$.

\punts{0,4} Esquema
\figura[0.22]{sol_fig3.png}

\punts{0,2} $F_{AC} = F_{BC}$

\punts{0,2} $\vec{F}_C = I_C\vec{l}_C\times\vec{B}_{total}$

\punts{0,2} $F = 0,30\times 2,0\times 4,2\times 10^{-7} = 2,55\times 10^{-7}\ \mathrm{N}$
\end{solapartat}
